\documentclass[11pt,a4paper]{article}

\usepackage[margin=2.5cm]{geometry}
\usepackage{amsmath, amssymb}
\usepackage{enumitem}
\usepackage{xcolor}
\usepackage{hyperref}

\usepackage[style=english]{csquotes}

\title{Questions}
\author{}
\date{\today}

\begin{document}

\maketitle

I have a few questions regarding the case study and the implementation details.

First, what \textbf{lens function}
\[
    f:\mathbf{R}^{28}\to\mathbf{R}^{2}
\]
was chosen in the case study? I found the description that the lens was chosen
\enquote{based on density and centrality}. Does this mean something of the
following form?
\[
    f(x) = (\rho(x),c(x)),
    \qquad x\in\mathbf{R}^{28},
\]
where
\[
    \rho(x)
    =
    \frac{1}{n}\sum_{j=1}^{n}
    K\left(\frac{d(x,x_j)}{h}\right)
    \qquad
    \mbox{(kernel density estimate)}
\]
and
\[
    c(x)
    =
    \sum_{j=1}^{n} d(x,x_j)
    \qquad
    \mbox{(centrality)}?
\]

Was the distance \(d\) in the lens computation the special block-wise Euclidean
distance
\[
    d(x_i,x_j)
    =
    \left\|x_i[1:14]-x_j[1:14]\right\|_2
    +
    \left\|x_i[15:28]-x_j[15:28]\right\|_2?
\]

Second, what kind of \textbf{cover}
\[
    {\cal U}=\{U_i\},
    \qquad U_i\subseteq\mathbf{R}^{2},
\]
was chosen? Is there an explicit formula for the cover, including the chosen
parameters?

Assume that we have applied the lens and the cover, with well-chosen parameters,
and have obtained a pullback cover
\[
    \{X_i\}.
\]
The new part is now that we apply our scale-selection principle to each \(X_i\)
separately. More precisely, for each \(X_i\), we consider the Vietoris--Rips
filtration
\[
    j\mapsto \mathrm{VR}_{\epsilon_j}(X_i),
\]
assemble the scale-similarity matrix
\[
    \Sigma^{(N)}
    =
    \left(
    \begin{array}{cccc}
        c_1 & c_1 & \cdots & c_1 \\
        c_1 & c_2 & \cdots & c_2 \\
        \vdots & \vdots & \ddots & \vdots \\
        c_1 & c_2 & \cdots & c_N
    \end{array}
    \right),
    \qquad
    c_j = F\left(\mathrm{VR}_{\epsilon_j}(X_i)\right),
\]
with the chosen functional
\[
    F(K)
    =
    \#\mbox{components}(K)\cdot \#\mbox{edges}(K).
\]
Then we choose the index
\[
    j_{\mathrm{crit}}
    \in
    \arg\max_j v_j^*,
\]
where \(v^*\) is the principal eigenvector of \(\Sigma^{(N)}\), and set
\[
    \epsilon_{\mathrm{crit}}
    =
    \epsilon_{j_{\mathrm{crit}}}.
\]
The local threshold is then
\[
    \alpha_i = \epsilon_{\mathrm{crit}}.
\]

Were the Vietoris--Rips complexes again constructed using the special block-wise
Euclidean distance?

Could you also briefly explain what was done before applying the spectral
scale-selection step?

What would you say are the main advantages of this new scale-selection procedure?
For example, would you describe it as simpler, automatic, data-dependent, and less
dependent on manual parameter tuning? Is it also computationally cheaper, or is
the main advantage primarily conceptual and methodological?

Finally, which clustering algorithm was used to obtain the four communities in the
end? Were these the same four communities as before, that is, before using the
spectral threshold choice, or were there noticeable differences?



\end{document}
